\documentclass[10pt,a4paper]{amsart}
\usepackage[margin=19mm]{geometry}
\usepackage{amsmath,amssymb,mathtools}
\usepackage[scr=rsfs,cal=euler]{mathalfa}
\usepackage{xfrac}
\usepackage[unicode,hidelinks,bookmarks=false]{hyperref}
\newcommand{\ie}{i.e.\ }
\newcommand{\cf}{cf.\ }
\newcommand{\resp}{respectively\ }
\newcommand{\Subsection}{\subsection}
\providecommand{\nobreakdash}{\nobreak\hbox{-}}
\setlength{\emergencystretch}{2em}
\multlinegap0pt
\usepackage{bm}
\usepackage{bbm}
\usepackage{dsfont}
\usepackage{xspace}
\usepackage{cancel}
\usepackage{mathtools}
\usepackage{amsthm}
\makeatletter
\@ifundefined{thm}{\newtheorem{thm}{Thm}}{}
\@ifundefined{lemma}{\newtheorem{lemma}[thm]{Lemma}}{}
\makeatother
\DeclareMathOperator{\inv}{inv}
\newcommand{\Br}{\operatorname{Br}}
\title[Positroid Links and Braid varieties]{Positroid Links and Braid varieties}
\author{Roger Casals, Eugene Gorsky, Mikhail Gorsky, Jos\'e Simental}
\date{Source: arXiv:2105.13948v4 (CC BY 4.0)}
\begin{document}
\maketitle

\noindent\emph{Source and licence.} Positroid Links and Braid varieties. Version arXiv:2105.13948v4 (CC BY 4.0), distributed under the Creative Commons Attribution 4.0 International licence, as recorded in the arXiv licence metadata for this exact version and shown on the abstract page https://arxiv.org/abs/2105.13948v4. Full rights notice: licenses/arxiv-2105.13948-CC-BY-4.0.txt.\par\smallskip

\noindent\textit{Editorial scope.} Counted unit: the Lemma whose source label is \texttt{lem: inductive destabilization} in the article (source0.tex lines 2682--2708, source label \texttt{lem: inductive destabilization}), delivered together with the article's own complete original proof of it (source0.tex lines 2710--2751, one \texttt{proof} environment of 42 lines). No mathematics is written, repaired or re-derived: statement and proof are the article's own text line for line. Required context retained uncounted: lem: inv (lines 1525--1532); eq: noninversions (lines 1542--1544); lem:Markovsmooth (lines 2526--2546). Every definition, display and label of those lines is retained unchanged. Omitted from this package and not redistributed: 1--1524, 1533--1541, 1545--2525, 2547--2681, 2709--2709, 2752--5176. Nothing inside a retained excerpt is omitted or altered: every retained body is delivered exactly as the article's lines are written, so no omission receipt of any kind accompanies this package. Citations: the retained excerpts carry no citation command at all, so no bibliography entry of the article is transcribed and no reference is redistributed. Reference adaptation: none was needed and none was made; every reference inside the delivered unit resolves inside it, so no unresolved reference or citation is produced. Printed slips kept literal: no printed word, symbol, hypothesis or number of the counted statement or of its proof was altered. Layout and class: the article's own class is \texttt{amsart} with packages including amsmath, amssymb, amsxtra, anysize, tikz, hyperref, tikz-cd, graphicx, combelow, comment, enumerate, enumitem, amsthm, amsfonts; the wrapper is the lane's pinned \texttt{amsart} editorial wrapper at 10pt on A4, which restates the theorem-like environments the retained text needs and adds the article's own macro definitions that the retained text needs (\texttt{\textbackslash Br}, \texttt{\textbackslash inv}), with extra wrapper packages (bm, bbm, dsfont, xspace, cancel, mathtools). The delivered numbering is the unit's own. Assets: no retained span contains a figure environment, a graphics-inclusion command, a PDF-page inclusion, a table or a TikZ picture; no image file is copied into this package, the package declares no asset dependency and no image file is redistributed with it. Licence: version arXiv:2105.13948v4 of this article is distributed under the Creative Commons Attribution 4.0 International licence, as recorded in the arXiv licence metadata for this exact version and shown on the abstract page https://arxiv.org/abs/2105.13948v4. A full rights notice is delivered at licenses/arxiv-2105.13948-CC-BY-4.0.txt. Characters: every retained excerpt is pure ASCII, so the delivered engine, XeLaTeX with its default Latin Modern text font, reports no missing character.
\begin{lemma}
\label{lem: inv}
For all $s\in [1,n-k]$, we have the inequalities
$$
j_s\le f(j_s)-\inv(s)\le k+s
$$
The first inequality is sharp unless $f(j_s)=j_s$.
\end{lemma}

\begin{equation}\label{eq: noninversions}
\sharp\{t>s: f(j_t)>f(j_s)\}=n-k-s-\inv(s),
\end{equation}

\begin{lemma}\label{lem:Markovsmooth} Let $\beta_1\in\Br_n$ be a positive braid on $n$ strands in the generators $\sigma_1,\ldots,\sigma_{n-1}$ and $\beta_2\in\Br_{n+1}$ the lift of a permutation $w\in S_{n+1}$ on $(n+1)$ strands. Set $j:=w^{-1}(n+1)$, let $\tilde{\beta}_1 \in \Br_{n+1}$ be the braid obtained from $\beta_1$
%\MG{from $\beta_1$, right?} 
by shifting the indices of the Artin generators up by $1$, and let $\beta'_2 \in \Br_n$ be the braid obtained from $\beta_2$ by removing the strand that ends at the bottom on the right of $\beta_2$.

%\MG{I moved the definition of $\beta'_2$ up, since it appears in both points (1) and (2). Commented it out in (1) and (2).}

\begin{itemize}
    \item[$(1)$] Consider the braid words $\eta_1 := \Delta_{n+1}\beta_2 (\sigma_{\ell}\cdot\ldots\cdot\sigma_{j+1}\sigma_j\cdot\ldots\cdot\sigma_2\sigma_1)\tilde{\beta}_1 \in \Br_{n+1}$, an $(n+1)$-stranded positive braid, and $\eta_2:=\Delta_n\beta'_2(\sigma_{\ell - 1}\cdot\ldots\cdot\sigma_{j+1})\beta_1 \in \Br_n$, an $n$-stranded positive braid.
    %, with $\beta'_2$ obtained from $\beta_2$ by removing the strand that ends at the bottom on the right of $\beta_2$.
    \\

\noindent  Then the 0-framed smooth closure of $\eta_1$ is smoothly isotopic to that of $\eta_2$.\\

\item[$(2)$] Consider the braid words $\eta_1:=\Delta_{n+1}\beta_2(\sigma_{j}\sigma_{j-1}\cdot\ldots\cdot\sigma_1)\tilde{\beta}_1 \in \Br_{n+1}$, an $(n+1)$-stranded positive braid, and $\eta_2:=\Delta_n\beta'_2\beta_1\in \Br_n$, an $n$-stranded positive braid.
%, with $\beta'_2$ obtained from $\beta_2$ by removing the strand that ends at the bottom on the right of $\beta_2$.
\\

\noindent  Then the 0-framed smooth closure of $\eta_1$ is smoothly isotopic to the unlinked union of the 0-framed smooth closure of  $\eta_2$ and an unknot.
\end{itemize}
\hfill$\Box$
\end{lemma}

\begin{lemma}
\label{lem: inductive destabilization}
At the $m$-th step, $1\le m\le n-k+1$, of the above destabilization procedure we get an $(n+1-m)$-strand braid  that can be decomposed into the following three parts:
\begin{itemize}
    \item[(a)] The first part $\beta^{(1)}_{m}$ is the product of interval braids
\[
\prod_{h = m}^{n-k}\left(s_{\lambda^{t}_{n-k+1-h}+h-1}\cdots s_h\right) \times
\prod_{\ell = 0}^{m-2}\left(s_{\lambda^{t}_{n-k-m-2+\ell}+n-k-1}\cdots s_{2n-k-m-\ell+3+\inv(n-k-m+2-\ell)-f(j_{n-k-m+2-\ell})}\right),
\]    
%\[
%\left(s_{\lambda_{n-k-m+1}^{t} + m -1}\cdots s_{m}\right)
%\cdots\left(s_{\lambda_{1}^{t} + n - k -1}\cdots s_{n-k}\right)
%\times   
%\] 
%\[
%\left(s_{\lambda^{t}_{n-k-m+2}+n-k-1}\cdots s_{2n+\inv(n-k-m+2)-f(j_{n-k-m+2})-k-m+3}\right)\cdots
%\left(s_{\lambda^{t}_{n-k}+n-k-1}\cdots s_{2n-k+
%-f(j_{n-k})+1}\right).
%\]
where for $m=1$ the second product is assumed to equal 1 and for $m = n-k+1$ the first product is assumed to equal 1.

\vspace{0.1cm}

\item[(b)] The second part $\beta^{(2)}_{m}$ is a permutation braid obtained as a part of the wiring diagram for $u^{-1}w_0$ connecting $1, \dots, n-m+1$ on the right with $u^{-1}(n),\ldots,u^{-1}(m)$ on the left.\\%\textcolor{red}{No relabeling yet!}\\

\item[(c)] The third part is $\Delta_{n-m+1}$.   \end{itemize}
\end{lemma}

\begin{proof}
We prove this by induction in $m$. For notational convenience, we refer to the left end of $\beta^{(1)}_{m}$ as $x=0$, to the right end of $\beta^{(1)}_{m}$ which coincides the left end of $\beta^{(2)}_{m}$ as $x=1$ and to the right end of $\beta^{(2)}_{m}$ as $x=2$ (see figures above). For $m=1$ we get
\begin{equation}\label{eq: w0ww0}
%w_0ww_0
\beta^{(1)}_{1}=w^*= \left(s_{\lambda^{t}_{n-k}}\cdots s_1\right)\cdots\left(s_{\lambda_{1}^{t} + n - k -1}\cdots s_{n-k}\right). 
\end{equation}
and $\beta^{(2)}_{1}=u^{-1}w_0$. 

Assume that the statement holds for all $1 \leq j \leq m$.
To prove the step of induction, we mark in blue the bottom-most strand at $x=2$. It is labeled by $n-m+1$ at $x=2$ and by $(u^{-1}w_0)^{-1}(n-m+1)  = w_0u(n-m+1) = w_0(f(j_{n-k-m+1})) = n+1-f(j_{n-k-m+1})$ at $x=1$. For brevity, denote $m' := n-k-m+1$.

By assumption, the rightmost interval braid in $\beta^{(1)}_{m}$ is
%ends with the interval braid
\begin{equation}\label{eq: m step}
s_{m-1+\lambda^{t}_{n-k-m+1}}\cdots s_{m}=s_{m-1+\lambda^{t}_{m'}}\cdots s_{m}.
\end{equation}

We need to verify that we can apply Lemma \ref{lem:Markovsmooth}. Let us count the number of strands \emph{above} the blue strand $n+1-f(j_{m'})$ we have removed at $x=1$. Note that such strands correspond to $1 \leq \ell < m$ such that $n+1-f(j_{\ell'}) < n+1-f(j_{m'})$ where $\ell' :=  n-k-\ell+1$, that is, they correspond to elements $n - k \geq \ell' > m'$ such that $f(j_{\ell'}) > f(j_{m'})$. Thanks to \eqref{eq: noninversions}, there are precisely $n-k-m'-\inv(m')$ such strands. We conclude that the blue strand occupies the position
\[
n + 1 - f(j_{m'}) - (n-k-m' - \inv(m')) = k+1+m'+\inv(m') - f(j_{m'})
\]
and, in order to verify we can apply Lemma \ref{lem:Markovsmooth}, we need to verify that 
\[
k+m'+\inv(m') - f(j_{m'}) \leq \lambda^{t}_{m'} = k+m'-j_{m'} \Leftrightarrow j_{m'} \leq f(j_{m'}) - \inv(m')
\]
which holds thanks to Lemma \ref{lem: inv}. So we can indeed use Lemma \ref{lem:Markovsmooth} in order to destabilize. We have cases:

\begin{itemize}
\item[-] If $j_{m'} = f(j_{m'}) - \inv(m')$, then we remove a disjoint blue unknot and erase the interval braid completely.
\item[-] If $j_{m'} = f(j_{m'}) - \inv(m') - 1$, then we also erase the interval braid completely.
\item[-] Else, we delete the part $(s_{m+k+m'+\inv(m')-f(j_{m'})}\cdots s_{m})$ intersecting the blue strands out of the interval braid \eqref{eq: m step} and are left with 
\[
s_{m-1+\lambda^{t}_{m'}}\cdots s_{m+k+m'+\inv(m')-f(j_{m'})+1} = s_{m-1+\lambda^{t}_{m'}}\cdots s_{n+\inv(m')-f(j_{m'}) + 2}.
\]
Now we push this interval through the $m' - 1$ remaining interval braids and get
\begin{equation}\label{eq: push before substract}
s_{m-1+\lambda^{t}_{m'}+m'-1}\cdots s_{n+\inv(m')-f(j_{m'}) + 2+m'-1}=
s_{n-k-1+\lambda^{t}_{m'}}\cdots s_{n+\inv(m')-f(j_{m'}) + m'+1}.
\end{equation}
This agrees with the interval braid in $\beta^{(1)}_{m+1}$.
\end{itemize}
\end{proof}

\end{document}
